\begin{lemma}
Let $g\in\mathsf{IncSeq}$ be nonidentity.  There is a prefix-continuous map
$\sigma:\mathsf{IncSeq}\to\mathsf{IncSeq}$ satisfying
\[
\sigma(f\circ\mathsf{SuccSeq})=\sigma(f)\circ g
\]
for every $f$.  Equivalently, the ordinary shift continuously
homomorphically maps to right composition by $g$.
\end{lemma}
\begin{proof}
Choose $k_g$ from \noderef{FirstMovedPoint} and put
$G(n)=g^{[n]}(k_g)$.  For each $f$, use the block formula of
\noderef{BlockSigma} to define $\sigma_f$.  The node
\noderef{BlockSigmaEmbedding} says that this is an increasing injection,
and \noderef{BlockSigmaContinuous} says that the assignment
$f\mapsto\sigma_f$ is prefix-continuous.  Thus it gives a candidate map
$\sigma$ into $\mathsf{IncSeq}$.

Finally, \noderef{BlockSigmaIntertwines} gives, coordinate by coordinate,
\[
\sigma_{f\circ\mathsf{SuccSeq}}(l)=\sigma_f(g(l)).
\]
Therefore $\sigma(f\circ\mathsf{SuccSeq})=\sigma(f)\circ g$, so the
candidate is the required continuous homomorphism.
\end{proof}
